\section{Conception}
\label{sec:conception}

\begin{resume}
La recherche par vagues s\'epare strictement deux temps. Pendant la collecte,
plusieurs workers descendent en parall\`ele avec leurs propres plateaux et
n'\'ecrivent jamais les statistiques de l'arbre ; ils r\'eservent les feuilles
par une transition atomique et d\'eposent des r\'esultats immuables. Apr\`es la
barri\`ere, le coordinateur ex\'ecute une seule inf\'erence, d\'eveloppe les
feuilles et remonte les valeurs, s\'equentiellement. La correction tient \`a
quatre m\'ecanismes : des \'etats de n\oe{}ud atomiques, une garde RAII qui
poss\`ede le chemin r\'eserv\'e, une table de transposition par snapshots
sous verrous ray\'es, et un budget de simulations exact. La production
continue, les verrous par n\oe{}ud et la d\'eduplication des \'evaluations
sont explicitement \'ecart\'es \`a ce stade.
\end{resume}

\subsection{Vue d'ensemble d'une vague}

Une vague est la plus petite unit\'e de travail du syst\`eme. Le coordinateur
choisit un nombre de permis \'egal \`a $\min(\text{lot}, \text{simulations
restantes})$, puis les workers actifs descendent l'arbre en parall\`ele. Un
worker peut encha\^iner plusieurs tentatives par vague, de sorte que deux ou
quatre workers peuvent remplir un lot de huit. Quand tous les collecteurs
sont arr\^et\'es, la barri\`ere s'ouvre : le coordinateur fusionne les
r\'esultats dans un ordre stable, traite les terminaux, lance au plus une
inf\'erence group\'ee, d\'eveloppe les feuilles et effectue les remont\'ees.
La vague suivante ne commence qu'apr\`es publication compl\`ete de ces
modifications.

\begin{figure}[t]
\centering
\begin{tikzpicture}[x=1cm,y=0.62cm,font=\scriptsize]
% axe du temps
\draw[-{Latex[length=2mm]}, gris] (0,0.2) -- (13,0.2)
  node[right, font=\scriptsize] {temps};
% lignes de vie
\foreach \y/\lab in {1.4/W1, 2.2/W2, 3.0/W3, 4.2/Coord.} {
  \draw[gris!40] (0,\y) -- (12.6,\y);
  \node[left, font=\scriptsize, align=right] at (0,\y) {\lab};
}
% collecte parallele
\foreach \x/\y in {0/1.4, 0/2.2, 0/3.0} {
  \fill[bleu!55, rounded corners=1pt]
    (\x+0.15,\y-0.18) rectangle (\x+4.2,\y+0.18);
}
\node at (2.2,1.9) [font=\scriptsize, align=center]
  {descentes parall\`eles\\\scriptsize (plateaux priv\'es, virtual loss)};
\draw[dashed, ambre, thick] (4.5,0.4) -- (4.5,4.6)
  node[above, font=\scriptsize, ambre] {barri\`ere};
% inference + traitement
\fill[vert!45, rounded corners=1pt] (4.8,4.02) rectangle (8.6,4.38);
\node at (6.7,4.2) {une inf\'erence group\'ee};
\fill[gris!25, rounded corners=1pt] (8.6,4.02) rectangle (12.3,4.38);
\node at (10.45,4.2) {expansion + remont\'ees};
% vague suivante
\foreach \x/\y in {9.2/1.4, 9.2/2.2, 9.2/3.0} {
  \fill[bleu!25, rounded corners=1pt]
    (\x,\y-0.18) rectangle (\x+3.4,\y+0.18);
}
\node at (10.9,1.9) [font=\scriptsize] {vague suivante};
\end{tikzpicture}
\caption{D\'eroulement d'une vague. La collecte est parall\`ele et sans
\'ecriture des statistiques ; la barri\`ere garantit que l'inf\'erence, les
expansions et les remont\'ees voient un arbre stable.}
\label{fig:vague}
\end{figure}

\subsection{Pourquoi un arbre partag\'e}

L'alternative naturelle, un arbre par c\oe{}ur, a \'et\'e \'ecart\'ee pour
trois raisons. D'abord elle duplique les \'evaluations : deux arbres
ind\'ependants red\'ecouvrent la m\^eme position et paient deux inf\'erences.
Ensuite elle dilue la profondeur : huit arbres \`a 700 simulations explorent
chacun un arbre de 700 n\oe{}uds au lieu d'un arbre de 5\,600. Enfin elle
complique fortement la r\'eutilisation de racine : apr\`es un coup jou\'e, le
sous-arbre du coup doit \^etre retrouv\'e dans chaque arbre, ou perdu.

Le self-play, lui, fonctionne d\'ej\`a par parties ind\'ependantes : ses huit
fils OpenMP ne partagent pas d'arbre, et c'est tr\`es efficace parce que les
parties divergent naturellement. La recherche UCI n'a pas cette chance : une
seule position \`a analyser, donc un seul arbre. C'est pr\'ecis\'ement ce cas
que les vagues traitent.

\subsection{Le r\'esultat de collecte}

Une descente produit un objet d\'epla\c{c}able, \code{LeafWork}, qui ne
contient aucun pointeur vers le plateau local. Ses champs :

\begin{itemize}[leftmargin=1.4em]
  \item le type de r\'esultat : \code{Network}, \code{Terminal} ou
  \code{Collision} ;
  \item le n\oe{}ud r\'eserv\'e et la garde \code{PathReservation} compl\`ete
  qui porte le virtual loss de tout le chemin ;
  \item la cl\'e de cache et les coups l\'egaux, captur\'es \`a la feuille ;
  \item le tenseur du r\'eseau, uniquement pour une vraie feuille ;
  \item la valeur terminale, lorsqu'elle est d\'ej\`a connue.
\end{itemize}

Le r\'esultat transf\`ere au coordinateur la responsabilit\'e de terminer ou
d'annuler la r\'eservation exactement une fois. Cette propri\'et\'e est ce qui
rend les chemins d'erreur simples : la destruction de la garde suffit \`a
tout nettoyer.

\subsection{\'Etats de n\oe{}ud et propri\'et\'e}

Chaque n\oe{}ud porte un \'etat atomique \`a quatre valeurs :

\begin{center}
\small
\begin{tabular}{lp{10.5cm}}
\toprule
\'Etat & Signification \\
\midrule
\code{Unexpanded} & n\oe{}ud jamais r\'eserv\'e, ses enfants n'existent pas \\
\code{Pending}    & une descente le poss\`ede ; ses enfants se construisent \\
\code{Expanded}   & liste d'enfants compl\`ete et publi\'ee \\
\code{Terminal}   & mat, pat, nulle de r\`egle ou mat\'eriel insuffisant \\
\bottomrule
\end{tabular}
\end{center}

La propri\'et\'e s'acquiert par une comparaison-\'echange
\code{Unexpanded -> Pending}. Un seul worker gagne ; les autres d\'eclarent
une collision. \`A aucun moment deux descentes ne peuvent d\'evelopper le
m\^eme n\oe{}ud, ce qui r\'esout le pi\`ege historique du batching : deux jeux
d'enfants construits en parall\`ele sur la m\^eme feuille. L'\'etat
\code{Expanded} n'est publi\'e qu'apr\`es un \code{swap} de la liste
d'enfants compl\`ete, avec une \'ecriture \code{release} ; un lecteur fait
d'abord un chargement \code{acquire} de l'\'etat, puis lit la liste. Aucun
lecteur ne peut donc observer une liste partiellement construite.

\begin{figure}[t]
\centering
\begin{tikzpicture}[
  font=\small, node distance=3.4cm,
  etat/.style={draw=bleu!60, fill=bleuclair, rounded corners=2pt,
    minimum width=2.3cm, minimum height=0.8cm, align=center},
  trans/.style={-{Latex[length=2mm]}, draw=gris, thick}
]
\node[etat] (unexp) {Unexpanded};
\node[etat, right of=unexp] (pending) {Pending};
\node[etat, right of=pending] (expanded) {Expanded};
\node[etat, below of=pending, yshift=-0.4cm] (terminal) {Terminal};

\draw[trans] (unexp) -- node[above, font=\scriptsize]
  {CAS (un seul gagnant)} (pending);
\draw[trans] (pending) -- node[above, font=\scriptsize, align=center]
  {publication apr\`es\\construction des enfants} (expanded);
\draw[trans] (pending) -- node[right, font=\scriptsize, align=left, xshift=2pt]
  {mat, pat,\\nulle de r\`egle} (terminal);
\draw[trans, dashed] (pending) to[bend left=38]
  node[below, font=\scriptsize] {collision : release} (unexp);
\draw[trans] (expanded) to[loop above] node[above, font=\scriptsize]
  {s\'election} (expanded);
\end{tikzpicture}
\caption{Machine \`a \'etats d'un n\oe{}ud. La transition vers \code{Expanded}
est une publication : les enfants sont enti\`erement construits avant que
l'\'etat ne devienne visible. La transition en pointill\'es est le retour
arri\`ere d'une collision ou d'une exception.}
\label{fig:etats}
\end{figure}

\subsection{La garde RAII et le virtual loss}

Le virtual loss doit \^etre visible pendant la s\'election, pas apr\`es
l'arriv\'ee \`a la feuille. Chaque worker ajoute donc la racine puis chaque
enfant choisi \`a sa garde \code{PathReservation} avant de continuer la
descente. Deux workers peuvent encore choisir le m\^eme enfant
simultan\'ement ; le virtual loss r\'eduit cette probabilit\'e mais ne la
supprime pas. La garantie de correction reste la transition exclusive vers
\code{Pending}.

La garde poss\`ede deux choses : le chemin r\'eserv\'e (un vecteur de
pointeurs) et, le cas \'ech\'eant, le n\oe{}ud en \code{Pending} dont elle est
propri\'etaire. Sa destruction lib\`ere tout exactement une fois.

\begin{lstlisting}[caption={La r\'eservation, coeur de la garde RAII. La
r\'eservation enregistre le pointeur avant d'incr\'ementer le compteur, de
sorte qu'une allocation rat\'ee ne perd aucun compteur. La publication
d\'esarme la propri\'et\'e avant de pouvoir annuler un \'etat publi\'e.},
label={lst:reservation}]
bool PathReservation::try_claim(MCTSNode* node) {
    if (node == nullptr || m_owned_pending != nullptr) return false;
    NodeState expected = NodeState::Unexpanded;
    if (!node->state.compare_exchange_strong(
            expected, NodeState::Pending,
            std::memory_order_acq_rel, std::memory_order_acquire)) {
        return false;   // collision : un autre collecteur possede le noeud
    }
    m_owned_pending = node;
    return true;
}

void PathReservation::publish(NodeState state) noexcept {
    if (m_owned_pending == nullptr) return;
    m_owned_pending->state.store(state, std::memory_order_release);
    m_owned_pending = nullptr;   // la garde ne peut plus annuler
}

void PathReservation::release() noexcept {
    if (m_owned_pending != nullptr) {
        NodeState expected = NodeState::Pending;
        m_owned_pending->state.compare_exchange_strong(
            expected, NodeState::Unexpanded,
            std::memory_order_release, std::memory_order_relaxed);
        m_owned_pending = nullptr;
    }
    for (MCTSNode* node : m_path) {
        node->n_in_flight.fetch_sub(1, std::memory_order_relaxed);
    }
    m_path.clear();
}
\end{lstlisting}

Trois propri\'et\'es de cette garde m\'eritent d'\^etre soulign\'ees. D'abord,
\code{reserve()} enregistre le pointeur dans le vecteur avant d'incr\'ementer
le compteur atomique : si l'allocation \'echoue, aucun compteur n'est perdu.
Ensuite, \code{release()} ne remet en \code{Unexpanded} que le n\oe{}ud dont
la garde est propri\'etaire et qui n'a pas \'et\'e publi\'e ; si un autre
worker a repris le n\oe{}ud entre-temps, le CAS \'echoue et rien n'est
\'ecras\'e. Enfin, la publication d\'esarme la garde : un \code{release()}
ult\'erieur ne peut plus annuler un \'etat d\'ej\`a visible.

\subsection{La table de transposition par snapshots}

La table ne rend plus jamais de pointeur vers une entr\'ee. Un
\code{probe()} copie la valeur et les seules entr\'ees de politique utiles
dans un snapshot local, sous le verrou d'une bande ; un \code{store()}
pr\'epare ses donn\'ees puis remplace l'entr\'ee sous le m\^eme verrou.
L'index est le hash modulo la taille, et la bande est d\'eriv\'ee de l'index
physique r\'eel, pas d'un second modulo ind\'ependant sur le hash :

\begin{lstlisting}[caption={Index et bande. La bande est d\'eriv\'ee de
l'index r\'eellement acc\'ed\'e, ce qui garantit que deux cl\'es qui
s'\'ecrasent dans la m\^eme entr\'ee prennent toujours la m\^eme bande.},
label={lst:stripe}]
std::size_t EvaluationCache::index_for(const EvaluationCacheKey& key) const {
    return key.position_hash % m_entries.size();
}
std::size_t EvaluationCache::stripe_for(std::size_t index) const {
    return index % m_stripes.size();   // 4096 bandes au plus
}
\end{lstlisting}

L'ordre des rejets est celui de la s\'emantique \code{h0} : \`a hash
\'egal, on v\'erifie le compteur des cinquante coups, puis le contexte
courant, puis l'historique. L'attention \`a cet ordre n'est pas
cosm\'etique : il d\'etermine la classification des refus dans les compteurs,
donc la lisibilit\'e des campagnes. Un snapshot pris avant un
\'ecrasement reste valide ; c'est test\'e explicitement.

\subsection{Le succ\`es de table pendant la descente}

Sur un succ\`es de table, le worker qui poss\`ede le \code{Pending} copie le
snapshot, construit les enfants localement, publie \code{Expanded} et
\emph{poursuit la m\^eme descente}. Ce comportement est celui du
\code{select\_leaf} historique, et il a \'et\'e conserv\'e volontairement :
un succ\`es de table ne consomme pas une simulation et ne provoque pas de
remont\'ee \`a lui seul. Les autres workers ne lisent la liste d'enfants
qu'apr\`es avoir observ\'e \code{Expanded} par une lecture \code{acquire}.

Les n\oe{}uds terminaux, eux, sont d\'etect\'es avant toute consultation de
la table. C'est une garantie de s\'emantique : une nulle par r\'ep\'etition
ou par r\`egle des cinquante coups d\'epend de l'historique complet, et ne
doit jamais \^etre servie par une entr\'ee de cache calcul\'ee ailleurs.

\subsection{Budget exact, permis et collisions}

Le budget est la contrainte la plus facile \`a violer en silence. Trois
r\`egles le garantissent.

\begin{enumerate}[leftmargin=1.4em]
  \item Chaque vague dispose de $\text{slots} = \min(\text{lot},
  \text{simulations restantes})$ permis, et un CAS d\'ecr\'emente un compteur
  \code{available} avant toute tentative. Un permis n'est consomm\'e que par
  un r\'esultat final (terminal ou feuille r\'eseau) ; une collision le rend.
  \item Un compteur global d'essais borne la vague \`a $4 \times \text{slots}
  + \text{workers}$ tentatives, pour qu'un arbre tr\`es collisionn\'e ne fasse
  pas boucler les workers ind\'efiniment.
  \item Si les collisions emp\^echent de remplir le lot, le coordinateur
  ex\'ecute un lot partiel : le progr\`es est toujours garanti, et une vague
  vide est trait\'ee comme une erreur de logique, pas ignor\'ee.
\end{enumerate}

\`A la fin d'un appel r\'eussi, le nombre de remont\'ees neuves est exactement
\'egal au nombre de simulations demand\'e. Les compteurs distinguent les
positions \'evalu\'ees, les appels physiques, les succ\`es de table et les
collisions, ce qui permet de v\'erifier cette \'egalit\'e sur chaque
campagne.

\subsection{Exceptions, annulation et quiescence}

Toute exception, qu'elle vienne d'un worker, de l'\'evaluateur ou du
coordinateur, suit le m\^eme chemin : publication de l'annulation, r\'eveil
de tous les workers, attente de leur retour, annihilation des gardes encore
actives, restauration des \code{Pending} non publi\'es, puis relance de
l'exception. Le destructeur de la garde fait le gros du travail ; le
coordinateur ne fait que d\'etruire les r\'esultats qu'il d\'etient encore.

Au repos, l'arbre doit \^etre quiescent : aucun \code{Pending}, aucun
\code{n\_in\_flight} non nul, chaque r\'eservation termin\'ee ou annul\'ee
exactement une fois. Ces invariants sont v\'erifi\'es par un parcours
d'observabilit\'e, \code{inspect\_tree}, qui d\'etaille aussi les enfants
dupliqu\'es, les parents incoh\'erents, les priors qui ne somment pas \`a 1 et
les statistiques non finies.

\subsection{Instrumentation et observabilit\'e}

Les phases mesur\'ees sont la s\'election, la pr\'eparation du tenseur et de
la cl\'e, la lecture et l'\'ecriture de table, l'attente de verrou, la copie
du plateau, l'assemblage du lot, l'inf\'erence, l'expansion, les remont\'ees
et l'attente des workers. Le chronom\'etrage est optionnel, d\'esactiv\'e par
d\'efaut, et chaque thread \'ecrit dans son propre accumulateur, fusionn\'e
apr\`es la recherche. Les d\'efinitions sont volontairement inclusives : le
rapport ne pr\'etend pas mesurer des noyaux GPU, il mesure des temps d'appel.

\subsection{Ce qui a \'et\'e \'ecart\'e, et pourquoi}

\begin{center}
\small
\begin{tabular}{p{4.2cm}p{10.5cm}}
\toprule
Piste \'ecart\'ee & Raison \\
\midrule
Production continue & Deuxi\`eme \'etape, \`a d\'ecider apr\`es mesure. Elle
exige une synchronisation des statistiques et de la structure d'arbre que la
barri\`ere rend inutile. \\
Un verrou par n\oe{}ud & Consommation m\'emoire et contention pour un
b\'en\'efice inf\'erieur \`a la transition CAS, qui n'utilise que 8 bits
d'\'etat. \\
Atomiques sur $N$ et $W$ & Le protocole de lecture coh\'erente de $Q$ aurait
\'et\'e plus complexe que la barri\`ere, sans gain mesurable. \\
D\'eduplication des \'evaluations & Deux n\oe{}uds logiques diff\'erents
peuvent demander la m\^eme position ; d\'edupliquer est un chantier
d'optimisation distinct, hors du plan. \\
R\'eglage de $c_{\text{puct}}$ & Changer la recherche pendant qu'on change
l'ordonnancement rendrait toute mesure ininterpr\'etable. \\
Softmax sur les seuls coups l\'egaux & Optimisation ind\'ependante, not\'ee au
backlog. \\
\bottomrule
\end{tabular}
\end{center}
